\documentclass[12pt,ngerman,pointlessnumbers, abstracton, headsepline]{scrreprt}
\renewcommand{\familydefault}{\sfdefault}
\usepackage[T1]{fontenc}
\usepackage[latin9]{inputenc}
\usepackage{geometry}
\geometry{verbose,tmargin=2.5cm,bmargin=2.5cm}
\pagestyle{plain}
\setlength{\parskip}{\medskipamount}
\setlength{\parindent}{0pt}
\usepackage{color}
\usepackage{babel}
\usepackage{array}
\usepackage{varioref}
\usepackage{textcomp}
\usepackage{multirow}
\usepackage{amsmath}
\usepackage{amsthm}
\usepackage{graphicx}
\usepackage{setspace}
\usepackage{nomencl}
% the following is useful when we have the old nomencl.sty package
\providecommand{\printnomenclature}{\printglossary}
\providecommand{\makenomenclature}{\makeglossary}
\makenomenclature
\setstretch{1.2}

\makeatletter

%% Because html converters don't know tabularnewline
\providecommand{\tabularnewline}{\\}

%%%%%%%%%%%%%%%%%%%%%%%%%%%%%% User specified LaTeX commands.
% verschieden Symbole, Zeichen wie (c), �
\usepackage{textcomp,units}

% Mehr Platz zwischen Tabelle und Untertitel
\usepackage{caption}
\captionsetup[table]{skip=10pt}

%Kapitelzahl sehr gro�
\makeatletter% siehe De-TeX-FAQ 
 \renewcommand*{\chapterformat}{% 
   \begingroup% damit \unitlength-�nderung lokal bleibt 
     \setlength{\unitlength}{1mm}% 
     \begin{picture}(10,10)(0,5) 
       \setlength{\fboxsep}{0pt} 
       %\put(0,0){\framebox(20,40){}}% 
       %\put(0,20){\makebox(20,20){\rule{20\unitlength}{20\unitlength}}}% 
       \put(10,15){\line(1,0){\dimexpr 
           \textwidth-20\unitlength\relax\@gobble}}% 
       \put(0,0){\makebox(10,20)[r]{% 
           \fontsize{28\unitlength}{28\unitlength}\selectfont\thechapter 
           \kern-.05em% Ziffer in der Zeichenzelle nach rechts schieben 
         }}% 
       \put(10,15){\makebox(\dimexpr 
           \textwidth-20\unitlength\relax\@gobble,\ht\strutbox\@gobble)[l]{% 
             \ \normalsize\color{black}\chapapp~\thechapter\autodot 
           }}% 
     \end{picture} % <-- Leerzeichen ist hier beabsichtigt! 
   \endgroup 
}

\usepackage{ %a4wide,
            ellipsis, fixltx2e, mparhack,   %Fehlerkorrektur f�r Marginalien
            booktabs, longtable             %sch�nere Tabellen
}  

\usepackage[automark]{scrpage2}
%\automark[chapter]{chapter}
\clearscrheadfoot
\ohead{\\\headmark}
\ihead{\includegraphics[scale=0.15]{logo.jpg}}%\pagemark}
\ofoot[\pagemark]{\pagemark}


%Kurzfassung und Abstract (englisch) auf eine Seite
\renewenvironment{abstract}{
    \@beginparpenalty\@lowpenalty
      \begin{center}
        \normalfont\sectfont\nobreak\abstractname
        \@endparpenalty\@M
      \end{center}
}{
    \par
}



% sch�nerer Blocksatz!!
\usepackage{microtype}

\usepackage{ifpdf} % part of the hyperref bundle
\ifpdf % if pdflatex is used

%set fonts for nicer pdf view
 \IfFileExists{lmodern.sty}{\usepackage{lmodern}}
  {\usepackage[scaled=0.92]{helvet}
    \usepackage{mathptmx}
    \usepackage{courier} }
\fi

 % the pages of the TOC are numbered roman
 % and a pdf-bookmark for the TOC is added
 \pagenumbering{roman}
 \let\myTOC\tableofcontents
 \renewcommand\tableofcontents{
   %\pdfbookmark[1]{Contents}{}
   \myTOC
   \clearpage
   \pagenumbering{arabic}}

%Bezeichungen anpassen
%Babelpaket mu� zuvor geladen werden
%\usepackage[ngerman]{babel}
\addto\captionsngerman{ 
\renewcommand{\figurename}{Abb.}% 
\renewcommand{\tablename}{Tab.}% 
\renewcommand{\abstractname}{Kurzfassung}
%\renewcommand{\nomname}{Abk�rzungen}
}

% Alle Querverweise und URLs als Link darstellen
% In der PDF-Ausgabe
 \usepackage[colorlinks=true, bookmarks, bookmarksnumbered, bookmarksopen, bookmarksopenlevel=1,
  linkcolor=black, citecolor=black, urlcolor=blue, filecolor=blue,
  pdfpagelayout=OneColumn, pdfnewwindow=true,
  pdfstartview=XYZ, plainpages=false, pdfpagelabels,
  pdfauthor={LyX Team}, pdftex,
  pdftitle={LyX's Figure, Table, Floats, Notes, and Boxes manual},
  pdfsubject={LyX-documentation about figures, tables, floats, notes, and boxes},
  pdfkeywords={LyX, Tables, Figures, Floats, Boxes, Notes}]{hyperref}

%mehr Platz zwischen �berschrift und Tabelle
\newcommand{\@ldtable}{}
\let\@ldtable\table
\renewcommand{\table}{ %
                 \setlength{\@tempdima}{\abovecaptionskip} %
                 \setlength{\abovecaptionskip}{\belowcaptionskip} %
                 \setlength{\belowcaptionskip}{\@tempdima} %
                 \@ldtable}

%In dieser Arbeit wird auf die Nomenklatur als Abk�rzungsverzeichnis verzichtet. Bei Wunsch wieder aktivieren.
%Nomenklatur als Abk�rzungsverzeichnis verwenden
%\renewcommand{\nomname}{Abk�rzungsverzeichnis}
%\renewcommand{\nomlabelwidth}{20mm}

%Nomenklatur als Glossar verwenden
%Nur Noetig wenn auch Glossar verwendet wird.
\renewcommand{\nomname}{Glossar}

%Farbe f�r Programmcode festlegen
\definecolor{lightgray}{rgb}{0.8,0.8,0.8}

\AtBeginDocument{
  \def\labelitemiii{\(\circ\)}
}

\makeatother

\usepackage{listings}
\renewcommand{\lstlistingname}{Listing}

\begin{document}
\noindent \titlepage

\noindent \begin{center}
\begin{tabular}{cc}
 & \multirow{0}{*}{\includegraphics[width=5cm]{images/hn_niederrhein_logo_big}}\tabularnewline
{\large{}Fachbereich Elektrotechnik und Informatik}\hspace{1.5cm} & \tabularnewline
 & \tabularnewline
 & \tabularnewline
 & \tabularnewline
\end{tabular}
\par\end{center}

\noindent \vspace{7cm}


\noindent \begin{flushleft}
\textbf{\Large{}Projektfach WS 16/17}
\par\end{flushleft}{\Large \par}

\noindent \begin{flushleft}
{\large{}SLAM - Simultaneous Localization and Mapping}
\par\end{flushleft}{\large \par}

\noindent \begin{flushleft}
{\Large{}\vspace{1.5cm}
}
\par\end{flushleft}{\Large \par}

\begin{tabular}{ll}
eingereicht von:\hspace{1cm} & Andre K�rten\tabularnewline
 & Matrikelnummer: 111111\tabularnewline
 & Sebastian Schilling\tabularnewline
 & Matrikelnummer: 222222\tabularnewline
 & \tabularnewline
betreut durch: & Prof. Dr.-Ing. Thomas Nitsche\tabularnewline
 & Hochschule Niederrhein\tabularnewline
 & \tabularnewline
 & \tabularnewline
 & \tabularnewline
\multicolumn{2}{l}{Krefeld, der \today}\tabularnewline
\end{tabular}

\newpage{}

~

\vspace{17.1mm}


\noindent \begin{flushleft}
\textbf{\huge{}Kurzfassung}
\par\end{flushleft}{\huge \par}

Dieser Bericht fasst die Arbeiten und Ergebnisse des Wahlprojektfaches
des Masterstudiengangs Informatik der Hochschule Niederrhein zusammen.
Ziel des Projektes ist die Erstellung einer Umgebungskarte eines Roboters in einer f�r ihn unbekannten Gegend. Mithilfe dieser generierten Karte, soll es dem Roboter m�glich sein zu navigieren. Die oben genannten Anforderung lassen sich durch das SLAM-Problem beschreiben. F�r das Projekt wurden verschiedenen Algorithmen ausprobiert und miteinander verglichen.


\newpage{}

\tableofcontents{}

\newpage{}



\pagenumbering{roman}
\setcounter{page}{4}

\listoffigures


\newpage{}

\listoftables


\newpage{}

\pagenumbering{arabic}


\chapter{Einf�hrung}


\section{Problemstellung}

Segway..


\section{Was ist SLAM?}

In der Robotik steht das Akronym SLAM f�r Simultanesous Localization
and Mapping. Dabei beschreibt SLAM das Problem eines Roboters, eine Karte von seiner unbekannten Umgebung zu erstellen und gleichzeitig mit dieser zu navigieren.


\section{Partikel Filter}


\chapter{Datenformate}


\chapter{Algorithmen}

\section{DPSlam}

\section{TinySlam}

\section{MRPT}

\section{Sonstige}





\chapter{Ausblick}






% Dieser Code ist noetig, da sonst die falsche Seitenzahl im Inhaltsverzeichnis angezeigt wird
\clearpage
\phantomsection
% Die folgende Zeile sorgt dafuer, dass der Glossar im Inhaltsverzeichnis angezeigt wird.
\addcontentsline{toc}{chapter}{Glossar}

\nomenclature{APT}{Advanced Packaging Tool (kurz. APT) ist ein Paketverwaltungssystem das z.B. bei Debian und Ubuntu eingesetzt wird. Das Paketverwaltungssystem stellt einen Gro�teil der f�r das Betriebssystem verf�gbaren Programme und Erweiterungen bereit. Diese k�nnen durch den Benutzer meist in der Befehlszeile oder �ber eine grafische Oberfl�che installiert und konfiguriert werden. }\nomenclature{API}{Application Programming Interface (kurz API) ist eine Programmierschnittstelle die eine Anbindung von Au�en an ein bestehendes System bietet. Die Schnittstelle wird vom Hersteller der Software zur Verf�gung gestellt um Dritten die M�glichkeit eines Addons oder einer Erweiterung zu geben. Die Schnittstelle definiert dabei ein Standardformat und Standardbefehle.}

\settowidth{\nomlabelwidth}{APT}
\printnomenclature{}

% Dieser Code ist noetig, da sonst die falsche Seitenzahl im Inhaltsverzeichnis angezeigt wird
\clearpage
\phantomsection



\bibliographystyle{bibtex-daten/unsrtdin}
\addcontentsline{toc}{chapter}{\bibname}\nocite{*}
\bibliography{bibtex-daten/bachelorarbeit-info}

\end{document}
